\uuid{YOz4}
\titre{Équations complexes}
\auteur{bodin ; hanani ; mzari}
\organisation{exo7}
\licence{CC BY-NC-SA 4.0}
\datecreate{2020-01-20}
\niveau{L1}
\module{Algèbre}
\chapitre{Nombres complexes}
\sousChapitre{Racine carrée, équation du second degré}
\difficulte{3}
\theme{Équations, équations complexes}
\exo7qcmid{lille-1.nombres-complexes.025}

\contenu{
\begin{question}
\qclassification{104.02,104.03,104.04}
\qtags{motcle=équations complexes}
On considère l'équation $(E) :  \, (z^2+1)^2+z^2=0, \, z\in \Cc$. L'ensemble des solutions de $(E)$ est :
\begin{answers}
    \good{ $\{ \pm \frac{1+\sqrt5}{2}i \, ,\,  \pm \frac{1-\sqrt5}{2}i\}$}

    \bad{$\{\pm \frac{1+\sqrt5}{2} \, , \,  \pm \frac{1-\sqrt5}{2}\}$}

    \bad{$\{\pm \frac{1+\sqrt3}{2}i \, , \,  \pm \frac{1-\sqrt3}{2}i\}$}

    \bad{ $\{\pm \frac{1+\sqrt3}{2} \, , \,  \pm \frac{1-\sqrt3}{2}\}$}

\end{answers}
\begin{explanations}
Remarquer que $(z^2+1)^2+z^2= (z^2+1)^2 - (iz)^2= (z^2-iz+1)(z^2+iz+1)$. On peut aussi poser $Z=z^2$ et se ramener à une équation du second degré.
\end{explanations}
\end{question}
}
